% ---
% titre: Comparer des fractions
% theme: NO
% chapitre: Nombres rationnels
% sous_chapitre: Représenter, comparer et encadrer
% niveau: [10e]
% type: EVAL
% tags: [c-fractions, c-comparer, p-question-TS]
% difficulte: 1
% corrige: true
% ---

\question[4]
Compare les fractions. Transforme-les pour obtenir le dénominateur indiqué, complète les numérateurs, puis place le signe $<$, $>$ ou $=$ dans le cercle.

\vspace{10pt}
\begin{enumerate}[leftmargin=*,align=parleft, label=\alph*), itemsep=25pt]
\item $\dfrac{3}{4}$ et $\dfrac{4}{5}$ \quad \quad $\Longrightarrow$ \quad\quad $\dfrac{\sol{\;\bl}{15}}{20}$ \quad \sol{\tikz[baseline=-0.6ex]\draw (0,0) circle (0.35cm);}{$<$} \quad $\dfrac{\sol{\;\bl}{16}}{20}$
\item $\dfrac{2}{7}$ et $\dfrac{1}{3}$ \quad \quad $\Longrightarrow$ \quad\quad $\dfrac{\sol{\;\bl}{6}}{21}$ \quad \sol{\tikz[baseline=-0.6ex]\draw (0,0) circle (0.35cm);}{$<$} \quad $\dfrac{\sol{\;\bl}{7}}{21}$
\item $\dfrac{2}{3}$ et $\dfrac{5}{8}$ \quad \quad $\Longrightarrow$ \quad \quad$\dfrac{\sol{\;\bl}{16}}{24}$ \quad \sol{\tikz[baseline=-0.6ex]\draw (0,0) circle (0.35cm);}{$>$} \quad $\dfrac{\sol{\;\bl}{15}}{24}$
\item $\dfrac{4}{6}$ et $\dfrac{6}{9}$ \quad \quad $\Longrightarrow$ \quad\quad $\dfrac{\sol{\;\bl}{12}}{18}$ \quad \sol{\tikz[baseline=-0.6ex]\draw (0,0) circle (0.35cm);}{$=$} \quad $\dfrac{\sol{\;\bl}{12}}{18}$
\end{enumerate}

\begin{solution}[0pt]
\textbf{1pt par exercice :} 0.5pt pour les deux numérateurs corrects, 0.5pt pour le signe (accepté si cohérent avec les numérateurs de l'élève).
\end{solution}